\chapter{作品概述}

\section{需要同时核对两方身份的消息}

设想一个机构内的定向消息系统。Alice准备向Bob发送一条任务指令，她希望只有Bob能够读取；
Bob收到消息时还希望确认，这条指令确实来自自己认可的Alice。这里包含两个同等重要的条件：
接收者必须正确，发送者也必须正确。

普通公钥加密主要解决第一个条件。Alice使用Bob的公钥加密，其他人无法读取，但“Bob只接受
Alice发送的消息”通常还要依靠证书、签名和上层协议另行处理。数字签名主要解决第二个条件，
却不负责限制谁能看到消息。身份基匹配加密（Identity-Based Matchmaking Encryption，IB-ME）
把两项判断放进同一次消息恢复过程。

记发送者身份为$s$，发送者指定的接收身份为$r$；接收密钥所属身份为$r'$，接收者期望的发送
身份为$s'$。解密成功的条件为
\[
  (r=r')\ \land\ (s=s').
\]
Alice发送给Bob的密文，只有“Bob的接收密钥”和“期望发送者Alice”同时出现时才能打开。
表~\ref{tab:identity-cases}列出四种组合。任意一方不匹配，系统均返回拒绝。

\begin{table}[H]
\centering\small
\caption{接收身份与期望发送身份的四种组合}
\label{tab:identity-cases}
\begin{tabular}{ccc}
\toprule
\rowcolor{black!8} 接收身份 & 期望发送身份 & 结果 \\
\midrule
匹配 & 匹配 & 恢复消息 \\
匹配 & 不匹配 & 拒绝 \\
不匹配 & 匹配 & 拒绝 \\
不匹配 & 不匹配 & 拒绝 \\
\bottomrule
\end{tabular}
\end{table}

\begin{figure}[H]
\centering
\begin{tikzpicture}[node distance=13mm and 13mm]
 \node[flow,text width=34mm] (sender) {Alice的发送密钥\\指定接收者Bob};
 \node[flow,text width=30mm,right=of sender] (ct) {加密消息\\经公开信道传输};
 \node[check,text width=42mm,right=of ct] (receiver) {Bob的接收密钥\\期望发送者Alice};
 \node[check,text width=42mm,below=of receiver] (result) {双方身份匹配：恢复原文\\任一身份错配：拒绝};
 \draw[flowarrow] (sender)--(ct);
 \draw[flowarrow] (ct)--(receiver);
 \draw[flowarrow] (receiver)--(result);
\end{tikzpicture}
\caption{身份匹配加密的基本使用方式}
\label{fig:identity-match}
\end{figure}

身份可以是邮箱、组织账号、设备序列号，也可以是业务系统生成的规范标识。可信发钥机构完成
现实身份核验并签发相应密钥；通信双方在发送和接收时直接填写身份。该方式适合机构内定向
投递、受控终端之间的指令传输、限定消息来源的设备管理，以及需要隐藏通信关系的消息服务。

\section{从理论构造到可运行系统}

Wang等提出了由两层匿名分层身份基加密（HIBE）和身份基签名（IBS）构造IB-ME的方法
~\cite{wang2024ibme}。HIBE把接收者放在第一层、发送者放在第二层，负责沿身份路径加密；IBS
认证发送者及其签署的协议消息。两种原语可以分别选择不同方案，因此天然适合形成通用组合。

论文构造给出了清晰的密码关系，落到软件时还要处理四个问题。

\textbf{原语之间的容量不同。} 有些HIBE只适合加密几十字节，而IBS签名可能达到数KiB、数MiB，
应用消息也没有固定长度。若直接把“签名＋消息”送入HIBE，大部分组合会立即超过明文容量。

\textbf{对象结构和参数不同。} 两个HIBE、六个IBS来自普通矩阵格、环格、模格和NTRU结构格，
它们使用不同的维度、模数、身份映射、采样器和序列化格式。公共协议需要在不读取后端内部
秘密的前提下完成装配。

\textbf{论文公式与实际资源之间有明显距离。} 一个在公式中写成矩阵平方根或高斯采样的步骤，
展开后可能需要数十GiB甚至更大的中间对象。十进制JSON表示便于调试，却会让包含大量整数的
密钥和签名成倍膨胀。

\textbf{单次运行与连续服务的成本不同。} 建立主密钥、为用户发钥、为固定身份对委派叶密钥，
都属于可复用工作。若每收一条消息都重新执行这些步骤，完整系统的延迟会被重复准备成本主导。

因此，本作品不只实现一条“能加密、能解密”的演示路径，而是把底层原语、组合协议、状态
管理、命令行和浏览器操作组织成同一套系统，并针对实测瓶颈进行优化。

\section{作品总体方案}

系统的输入是HIBE选择、IBS选择、接收身份、发送身份和应用消息。装配层首先检查两项原语的
身份长度、消息容量、签名上限和编码格式，再建立完整运行实例。协议层始终调用统一接口，
底层方案的矩阵、多项式和陷门只在各自后端中出现。

发送时，系统随机生成32字节种子，用两层HIBE沿$(r,s)$路径封装该种子；再从种子导出对称
密钥，用AES-256-GCM加密IBS签名和应用消息。为了把发送身份与完整传输字节连接起来，协议将
一次性验签公钥的承诺写入IBS签名，并用对应的一次性私钥认证HIBE胶囊、对称密文和公开上下文。

接收时，系统先验证最外层字节认证，再由接收密钥和期望发送身份产生HIBE叶密钥。种子恢复后，
系统检查GCM标签、解析内部载荷、验证IBS签名，并核对内外两处一次性公钥承诺。只有全部步骤
通过，应用消息才会交给用户。

\begin{figure}[H]
\centering
\begin{tikzpicture}[node distance=9mm and 10mm]
 \node[flow,text width=34mm] (primitive) {格HIBE与格IBS\\独立后端};
 \node[flow,right=of primitive,text width=34mm] (assembly) {能力与容量检查\\组合配置};
 \node[flow,right=of assembly,text width=36mm] (protocol) {IB-ME五算法\\Hybrid与字节认证};
 \node[check,below=of assembly,text width=34mm] (state) {规范对象\\发钥与状态恢复};
 \node[check,right=of state,text width=36mm] (entry) {命令行与浏览器\\消息收发};
 \draw[flowarrow] (primitive)--(assembly);
 \draw[flowarrow] (assembly)--(protocol);
 \draw[flowarrow] (assembly)--(state);
 \draw[flowarrow] (protocol)--(entry);
 \draw[flowarrow] (state)--(entry);
\end{tikzpicture}
\caption{从底层格原语到应用入口的总体关系}
\label{fig:work-overview}
\end{figure}

常规组合矩阵包含LATTE、Lai两种HIBE，以及BNN＋ML-DSA、BNN＋Falcon、Luo、SRGG、Sageloli
五种IBS，共十种配对。它们共享相同的身份顺序、混合加密和传输格式。目标规模Wang--PW IBS
随后与LATTE完成独立组合，用于补充矩阵SIS路线。

\section{主要工作与创新点}

作品的主要工作可以概括为通用设计、原语实现、工程优化和完整系统四个部分。

\begin{table}[H]
\centering\small
\caption{作品的主要工作及其作用}
\label{tab:main-contributions}
\begin{tabularx}{\textwidth}{@{}p{.2\textwidth}p{.37\textwidth}X@{}}
\toprule
\rowcolor{black!8} 工作 & 具体内容 & 解决的问题 \\
\midrule
通用IB-ME设计 & 独立HIBE／IBS接口、五算法、32字节种子封装、完整字节认证、规范对象 & 让不同格结构和容量的原语进入同一协议 \\
多算法实现 & 两种HIBE、五种常规IBS的十种配对；Wang--PW目标档与LATTE组合 & 展示普通矩阵格、环格、模格和NTRU结构方案 \\
关键算法优化 & Sageloli结构化采样、SRGG范围感知编码、Lai接收会话 & 降低矩阵展开、对象存储和重复委派成本 \\
软件与实验 & 固定发钥、重启恢复、CLI、浏览器演示、阶段计时、尺寸与内存统计 & 把原语实现变成可操作、可测量的完整消息系统 \\
\bottomrule
\end{tabularx}
\end{table}

\textbf{第一，建立可替换的抗量子IB-ME组合框架。} HIBE和IBS使用彼此独立的参数对象与版本号，
新增后端只需实现相应接口、对象编码和能力说明。公共五算法、身份顺序与传输格式保持不变。

\textbf{第二，使大规模环格IBS进入实际运行。} Sageloli论文参数中的采样若直接展开为稠密矩阵，
内存和分解成本难以承受。实现保留多项式结构，并利用每个频率上的低秩形式计算所需平方根，
完成PARAM-I和PARAM-II的完整生命周期。

\textbf{第三，按数学范围编码大整数对象。} SRGG四类对象中的字段具有不同取值范围。实现分别采用
模$q$定长编码、带符号短整数编码和二元编码，并在解码时恢复原对象。三档主状态缩小约
5.05至5.46倍，签名缩小约2.55至2.65倍。

\textbf{第四，把高成本委派从逐消息路径移到会话准备。} 对固定接收者和发送者，Lai HIBE的叶密钥
可以在会话开始时生成并保存。后续每条消息仍独立验证和解密，但无需重复执行同一身份路径的
委派。已有实验中，会话准备约510秒，后续解密中位数为28.6秒。

\section{报告内容安排}

第~2章介绍HIBE、IBS、格困难问题和方案选择依据。第~3章给出IB-ME的五个算法、混合加密、
完整密文字节认证和代表组合。第~4章展开LATTE、Lai、BNN、SRGG、Sageloli与Wang--PW的
实现，并说明三项主要优化。第~5章介绍软件分层、后端接入、状态管理和实际操作流程。第~6章
汇总时间、尺寸、内存、错误输入和恢复测试。第~7章总结作品特点及后续工作。
